% BHU PYQ -- Transcribed & Corrected 100% Accurate
% Paper: ESSE11: Smartphone Geosciences
% Exam: B.Sc. (Hons.) Semester I Examination 2024-25 (NEP)
% Category: Skill Enhancement Course (Earth Science)
% Ref Code: CECONF/N/2024-25/4

\documentclass[11pt,a4paper]{article}
\usepackage[a4paper,top=2.5cm,bottom=2.5cm,left=2.5cm,right=2.5cm,headheight=14pt]{geometry}
\usepackage{amsmath,amssymb}
\usepackage[T1]{fontenc}
\usepackage[utf8]{inputenc}
\usepackage{lmodern,microtype,fancyhdr,enumitem,hyperref,lastpage,parskip}

\hypersetup{colorlinks=true,urlcolor=black,linkcolor=black,pdftitle={ESSE11: Smartphone Geosciences -- BHU Sem I 2024-25 (NEP SEC)}}
\pagestyle{fancy}\fancyhf{}
\fancyhead[L]{\small\textit{Banaras Hindu University}}
\fancyhead[R]{\small\textit{ESSE11: Smartphone Geosciences (SEC)}}
\fancyfoot[C]{\small Page \thepage\ of \pageref{LastPage}}

\newlist{parts}{enumerate}{2}
\setlist[parts,1]{label=(\alph*),leftmargin=*,itemsep=4pt,topsep=2pt}
\setlist[parts,2]{label=(\roman*),leftmargin=*,itemsep=3pt,topsep=2pt}
\newcommand{\pts}[1]{\hfill\textit{[#1]}}

\begin{document}

\begin{center}
  {\large\bfseries BANARAS HINDU UNIVERSITY}\\[2pt]
  {\normalsize B.Sc. (Hons.) Semester I Examination 2024-25 (NEP)}\\[6pt]
  \rule{0.8\linewidth}{0.5pt}\\[6pt]
  {\large\bfseries SKILL ENHANCEMENT COURSE (EARTH SCIENCE / GEOLOGY)}\\[4pt]
  {\large\bfseries Paper Code: ESSE11 --- Smartphone Geosciences}\\[4pt]
  {\normalsize Time: 2:30 Hours\hfill Maximum Marks: 70}\\[2pt]
  {\small Paper Code Ref: CECONF/N/2024-25/4}
\end{center}

\rule{\linewidth}{0.4pt}\smallskip
\noindent\textit{Instruction: Write your Roll No. at the top immediately on the receipt of this question paper.}\\[2pt]
\noindent\textit{Note: Answer five questions in all, selecting at least two questions each from Section-A and Section-B. Section-C is compulsory. All questions carry equal marks (14 marks each).}
\bigskip

\begin{center}
  \textbf{\large SECTION --- A}
\end{center}
\textit{(Answer at least two questions from this section)}
\bigskip

\noindent\textbf{Question 1}\pts{14} \\
Write brief details on the applications of various parts of a smartphone with the help of a schematic diagram.

\medskip\rule{\linewidth}{0.2pt}\medskip

\noindent\textbf{Question 2}\pts{14} \\
Write the definition of ten important branches of geosciences with respect to applications and uses of smartphones.

\medskip\rule{\linewidth}{0.2pt}\medskip

\noindent\textbf{Question 3} \textit{Write short notes on the following:}\pts{$4 \times 3.5 = 14$}
\begin{parts}
  \item Importance of ChatGPT
  \item Types of Artificial Intelligence (AI)
  \item Internet of Things (IoT)
  \item Google Earth
\end{parts}

\bigskip
\begin{center}
  \textbf{\large SECTION --- B}
\end{center}
\textit{(Answer at least two questions from this section)}
\bigskip

\noindent\textbf{Question 4}\pts{14} \\
Briefly describe ten important uses of multimedia tools in education.

\medskip\rule{\linewidth}{0.2pt}\medskip

\noindent\textbf{Question 5}\pts{14} \\
Write the names of ten important sensors used in smartphones. Briefly describe the importance of smartphone sensors in education.

\medskip\rule{\linewidth}{0.2pt}\medskip

\noindent\textbf{Question 6} \textit{Write in brief on any two of the following:}\pts{$2 \times 7 = 14$}
\begin{parts}[label=(\roman*)]
  \item Global Positioning System (GPS).
  \item Application of social-media apps in education.
  \item Digital compass app.
\end{parts}

\bigskip
\begin{center}
  \textbf{\large SECTION --- C}
\end{center}
\textit{(Compulsory)}
\bigskip

\noindent\textbf{Question 7} \textit{Answer the following:}\pts{$7 \times 2 = 14$}
\begin{parts}[label=(\roman*)]
  \item Who is the developer of Android Operating System?
  \item Write the names of two system apps used in smartphones.
  \item Write the names of two navigation apps.
  \item Write the names of two teaching and learning apps.
  \item Write the names of two shopping apps.
  \item Write the names of two Geology apps.
  \item Write two important uses of social media apps.
\end{parts}

\vfill\rule{\linewidth}{0.4pt}
\begin{center}\small
\href{https://bhu-exam.vercel.app/}{Click here to visit our website}\\[4pt]
\href{https://me-aryan.vercel.app/}{Click here to visit our contributor}
\end{center}

\end{document}
